\section{Introduction}

The first primitive shared-clock collisions are classified by
\begin{equation}
 7p+3r=16Q,\qquad p<r,
 \label{eq:resonance}
\end{equation}
with two sign-symmetric shared coordinates \cite{tpc226}.  A source-native polarization
compiler subsequently identified the contribution of such a collision as a bilinear
$\beta$--$w$ block \cite{tpc228}.  The remaining finite geometry could in principle
form a complicated graph, creating global spectral or combinatorial losses.

We show that this concern disappears completely.  Equation \eqref{eq:resonance}
separates its two endpoint types at $8Q/5$, and each endpoint determines its unique
partner.  The resonance graph is therefore a matching for every $Q$.  Its spectrum and
all symmetric-source energy questions reduce to independent two-prime blocks.

This yields a sharp local saving criterion.  It also exposes the next obstruction:
geometry does not tell us how much physical source mass lies on matched vertices, nor
whether that mass favors antisymmetric modes.  Those are arithmetic questions rather
than graph-theoretic ones.
